%%%%%%%%%%%%%%%%%%%%%%%%%%%%%%%%%
%
%       EXERCÍCIO
%
%%%%%%%%%%%%%%%%%%%%%%%%%%%%%%%%%

\definevalorquestao

\begin{exercicioBanco}[\valorquestao]
A pureza (em percentual mássico) de um reagente químico produzido em um reator contínuo de síntese pode ser modelada por uma distribuição normal com média de $98{,}0\%$ e um desvio padrão de $0{,}5\%$.
%
\begin{enumerate}[(a)]
    \item Qual a probabilidade de um lote produzido apresentar pureza menor que $98{,}8\%$?
    %
    \item Qual a probabilidade de a pureza de um lote estar entre $97{,}0\%$ e $97{,}5\%$?
    %
    \item Que nível de pureza é excedido por $97{,}5\%$ dos lotes produzidos?
\end{enumerate}
%
\end{exercicioBanco}

